\section{Numerical verification of the constants}
\label{sec:numerics}

The local laws predict constants, not only exponents.  Assertion-based scripts
therefore test closed-form instances at high precision; the full tables and
execution details are in the technical supplement.  For the fixed
\(2\times2\) family
\(G=\bigl(\begin{smallmatrix}0&b_1\\b_2&c\end{smallmatrix}\bigr)\),
\(\Theta=e_1e_1^\top\), the three basic normalized limits are:
\begin{center}
\begin{tabular}{@{}lll@{}}
\toprule
Regime & predicted & computed at the smallest \(\delta\)\\
\midrule
regular & \(-16/9=-1.777777778\) & \(-1.777777777\)\\
strict corank one & \(-5/(4\sqrt{15})=-0.3227486122\) & \(-0.3227486102\)\\
boundary & \(2^{1/3}=1.259921050\) & \(1.259921042\)\\
\bottomrule
\end{tabular}
\end{center}
On the same boundary instance, standard Huber instead gives
\((1-\lambda_\delta^{\rm hub})/\delta=1.99999997\) at
\(\delta=10^{-8}\), versus the prediction \(1/\beta=2\) in
\eqref{eq:huber-boundary-rate-main}.
For the admissible profile
\(\psi'(x)=x/(1+x^3)^{1/3}\) on \(x\ge0\), the predicted \(p=3\)
constant is \(1.2778862085\); the computed normalized displacement is
\(1.2778860473\) at \(\delta=10^{-9}\).  This checks that the exponent
changes with the tail, not merely between the two named profiles.

For a corank-two check, take
\(A^\star=\diag(1,1,0,0)\),
\(\Theta=\diag(3/8,3/8,1,1/2)\), and \(\lambda^\star=1\).
Here \(\eta_R=0.715080068980\) and
\(\|K_F-K_R\|_F=0.040996138377\).  At \(\delta=10^{-4}\), the
pseudo-Huber slope is \(-0.7150801\), while Huber gives the water-filling
slope \(-0.6000000000\) and recovers \(\Phi_F^\star\) to \(10^{-11}\).
Because \(P_{\rm big}(s)\) and \(\widetilde B(s)\) are constant on this
diagonal line, \(\partial_s\mathcal W=0\) proves
\(t_1=\mathcal M_R=0\).  A non-diagonal rational instance instead predicts
\(\|\mathcal M_R\|_F=0.671775381\), versus the scaled value
\(0.671776068\) at \(\delta=10^{-6}\).

The compact-root instance
\(A^\star=\diag(2,1,0)\), \(\Theta=\diag(1,-1,1)\) verifies the predicted
cubic coefficient \(-3/8\) and quadratic direction norm \(\sqrt{26}/8\).
After the second cancellation, the rational \(4\times4\) instance in the
supplement has \(\chi_2=1/3\) and converges monotonically to the quintic
coefficient \(-1/8\); its normalized quartic direction converges to the full
matrix in \eqref{eq:compact-root-quartic-main}.

The moving-family assertions cover boundary, critical, strict-side, and
oscillatory paths and verify both the cubic root and direction constants.
Complete data, matrices, and finite-scale sequences are in the supplement;
the checks test constants but imply no occurrence-frequency law.
